% =====================================================================
\providecommand{\figpath}{figures/}
% problem2_section.tex —— 问题二论文片段(可直接插入中文竞赛论文)
% 使用: 竞赛模板中 \usepackage{amsmath,amssymb,graphicx,booktabs,caption}
%       图文件位于 figures/ 目录, 路径按需调整
% 数据来源: 全部由 results/*.csv 与本目录 MATLAB 程序实测生成, 见 README.md
% =====================================================================
\providecommand{\figpath}{figures/}

\section{问题二的模型建立与求解：第二个检测点的选择策略}

\subsection{问题二的目标与局部坐标系}

问题二要求：已知某全向干扰源在第一个检测点 $S_1$ 处测得的示向度，给出第二个检测点 $S_2$ 的选择策略，以获得关于该干扰源较好的定位效果，并给出第二个检测点的候选区域。

我们建立以 $S_1$ 为原点、第一次\emph{实际示向度}方向为 $x$ 轴正方向的局部坐标系，单位取米。这是坐标变换而非假设 $S_1$ 位于目标圆域圆心：若给出 $S_1$ 的全局坐标与第一次示向度 $\theta_1$（全局方位角），任意局部坐标 $G$ 与全局坐标 $\tilde G$ 的转换为
\begin{equation}
  \tilde G = S_1 + R(\theta_1)\,G,\qquad
  R(\theta)=\begin{pmatrix}\cos\theta & -\sin\theta\\ \sin\theta & \cos\theta\end{pmatrix}.
\end{equation}
在局部坐标系下，第一次示向度恒为 $0$，测向误差满足 $|\varepsilon|\le \delta$，其中 $\delta=1^\circ=\pi/180$。记干扰源位置为 $G=(u,v)$，第二检测点为 $P=(x,y)$，第二次实际示向度为 $\beta$（单位弧度，按圆周意义处理 $0/360^\circ$ 跨越）。

默认情形（本文全部数值实验的基准）：目标圆域在局部坐标下不裁切首次示向度扇形，即源可行域取“完整标准扇区”（§\ref{sec:DE}）。目标圆域参数在程序中作为可选输入保留（\texttt{crop\_enabled} 等参数），启用后所有几何量再与目标圆相交，且候选域退化为默认 $E$ 的安全子区域。

\subsection{源可行域 $D$ 与保守安全候选域 $E$}\label{sec:DE}

由题意：源的有效接收半径 $R$ 未知但满足 $1000\le R\le 1500$；$S_1$ 处已获得正常示向度，故源距 $S_1$ 在 $(5,1500]$ 内且 $|\arg G|\le\delta$；距检测点不超过 $5$ 米时有信号则无法获得正常示向度（转入光学定位清除分支）。因此标准情形的源可行域为
\begin{equation}
  D=\bigl\{G=(u,v):\ u\ge0,\ -u\tan\delta\le v\le u\tan\delta,\ 5<\sqrt{u^2+v^2}\le1500\bigr\}.
\end{equation}

要保证第二站 \emph{必然} 获得信号，需 $P$ 到任意可能源位置的距离都不超过有效接收半径的下界 $1000$，即
\begin{equation}
  E=\bigcap_{G\in D} B(G,1000)
   =\bigl\{P:\ \|P-G\|\le1000,\ \forall G\in D\bigr\}.
\end{equation}
这是未利用“首次已成功接收”这一信息的保守选择（保守：$E$ 中的点对 $R\in[1000,1500]$ 的任何源都保证可测；但它不一定是所有可能安全策略中的最大区域）。

对未被目标圆域裁切的标准 $D$，$E$ 可精确化简为四个圆盘之交：
\begin{equation}\label{eq:fourdisks}
  E=\bigcap_{\substack{r\in\{5,1500\}\\ s\in\{-1,1\}}}
    \bigl\{(x,y):\ (x-r\cos\delta)^2+(y-s\,r\sin\delta)^2\le 1000^2\bigr\}.
\end{equation}
其依据是：固定 $P$ 后 $\|G-P\|$ 在 $D$ 上的最大值可逐射线讨论——沿每条从原点出发的射线，$\|G-P\|$ 关于距离 $r$ 是凸函数，故最大值在射线两端 $r\in\{5,1500\}$ 处取得；而半径为 $1500$（或 $5$）的圆弧上，$\|G-P\|$ 的最大值在两端点 $\theta=\pm\delta$ 处取得（圆弧张角 $2\delta=2^\circ<\pi$，$\cos$ 在该小弧上单峰）。于是最坏距离只需检查四个角点 $(r\cos\delta,\ s\,r\sin\delta)$，$r\in\{5,1500\},s\in\{-1,1\}$。内侧半径 $5$ 米端点取闭包不改变最坏接收距离约束。\emph{注意：这一有限约束化简依赖“扇形未被裁切、圆弧张角小于 $\pi$”等条件，并非任意扇形都只需检查角点；若目标圆域裁切了 $D$，式~\eqref{eq:fourdisks} 给出的仍是安全子区域（可能更保守），本文第一版即采用之。}若保守地把首站附近 $5$ 米内也纳入源域（不扣近源孔），则只需三个圆盘：$x^2+y^2\le1000^2$ 及两个以 $(1500\cos\delta,\pm1500\sin\delta)$ 为心、半径 $1000$ 的圆盘；本文采用更精细的四圆盘刻画。数值验证（对 $D$ 边界 $2400$ 点采样的暴力最坏距离比对）与四圆盘判据完全一致。

\subsection{定位精度的刻画：最小包围圆}\label{sec:mec}

固定 $P$，第二次实际读数为 $\beta$ 时，正常观测对应的定位区域为
\begin{equation}
  J(P,\beta)=D_{\mathrm{ang}}(P)\cap W(P;\beta,\delta),
  \qquad D_{\mathrm{ang}}(P)=\{G\in D:\ \|G-P\|>5\},
\end{equation}
其中 $W(P;\beta,\delta)$ 是以 $P$ 为顶点、中心方向 $\beta$、半角 $\delta$ 的前向楔形，用两条半平面不等式实现：
\begin{equation}\label{eq:wedge}
  \sin(\beta-\delta)(u-x)-\cos(\beta-\delta)(v-y)\le0,\qquad
  -\sin(\beta+\delta)(u-x)+\cos(\beta+\delta)(v-y)\le0 .
\end{equation}
$J(P,\beta)$ 的边界由测向射线、扇区边界线与圆弧组成，不一定是四边形。若某读数 $\beta$ 对应的 $J$ 为空，说明该读数不可能由正常观测产生，应跳过，不能计为一次“精确定位成功”。

以最小包围圆（minimum enclosing circle, MEC）刻画区域大小：
\begin{equation}
  \rho(P,\beta)=\min_{c}\sup_{G\in J(P,\beta)}\|G-c\|,
\end{equation}
$\rho$ 为最小包围圆半径，$c$ 为对应圆心。真实主评价函数取所有可能正常读数的最坏情形：
\begin{equation}\label{eq:f1}
  f_1(P)=\sup_{\beta\in\Theta(P)}\rho(P,\beta),\qquad
  \Theta(P)=\bigcup_{G\in D_{\mathrm{ang}}(P)}
    \bigl[\arg(G-P)-\delta,\ \arg(G-P)+\delta\bigr],
\end{equation}
即同时考虑所有可能源位置与所有允许的第二次测角误差（不能只用真实方向、零误差或随机误差代替）。实际观测后，仅当 $\rho(P,\beta)\le 20$ 时，包围圆圆心才是保证成功的光学检测位置（光学定位有效距离 $20$ 米）；问题二只要求较好的定位效果，不要求第二次检测后必然清除。

数值实现：对给定 $\beta$，$J(P,\beta)$ 为“直线与圆弧围成的凸域减去两个半径 $5$ 米小孔”。圆弧采用内接/外切正 $N$ 边形（$N=2048$）做内外近似，半平面用 Sutherland--Hodgman 裁剪；近源孔（原点孔与 $P$ 孔）在\emph{内近似}中减去外接多边形（半径 $5+5\times10^{-4}$，结果仍是子集），在\emph{外近似}中减去内接多边形（半径 $5-5\times10^{-4}$，结果仍是超集；若外近似不扣孔，当 $P$ 位于扇区内部时定位区域贴着 $P$ 孔，上界将无法收敛，这是实现中发现并修正的关键细节）。凸多边形区域的最小包围圆等于其顶点集的最小包围圆，故对顶点集用 Welzl 随机增量算法（基础 MATLAB 实现，无工具箱依赖，单点/两点/共线/重复/近共线均有专门处理）求得圆心与半径，集合包含关系给出
\begin{equation}
  \rho(J_{\mathrm{in}})\le\rho(J(P,\beta))\le\rho(J_{\mathrm{out}}).
\end{equation}
这里的上下界是普通浮点实现配合显式集合包含关系的数值界限，不是区间算术的形式化认证；浮点容差（$10^{-9}$ 量级）与上述数学界限在程序与论文中明确区分。

\subsection{两次检测不能普遍保证 20 米覆盖}\label{sec:no20}

\emph{定理（下界）：}对标准完整扇区，任意 $P\in E$ 都有
\begin{equation}\label{eq:lower257}
  f_1(P)\ge r_*=\frac{1500\sin\delta}{1+\sin\delta}\approx25.73\ \text{米}.
\end{equation}
\emph{证明要点：}扇区内部包含圆盘 $C=\bigl(1500/(1+\sin\delta),\,0\bigr)$、半径 $r_*$（该圆盘与两条扇区边界线及外弧均相切，故 $C$ 圆盘 $\subset D$）。由 $E$ 的四圆盘刻画，$E\subset\{P:\|P-(5\cos\delta,0)\|\le1000\}$，故 $\|P\|\le1005<1500/(1+\sin\delta)-r_*\approx1448$，即 $P$ 必在 $C$ 圆盘之外。取实际读数 $\beta$ 指向 $C$：射线 $P\to C$ 穿过该圆盘的一条完整直径（长度 $2r_*\approx51.46$ 米），直径上所有点方向相同、全部属于 $D$ 且距 $P$ 超过 $5$ 米，因此相应定位区域包含这条直径，覆盖半径不小于 $r_*$。证毕。

由此，$f_1(P)\ge25.73>20$ 对所有 $P\in E$ 成立：\emph{两次检测不能保证把定位区域压缩到 20 米光学清除门槛之内}，问题二的目标只能是尽量降低 $f_1$。注意式~\eqref{eq:lower257} 是理论下界，不是已经求出的最优值，不能把 $25.73$ 米当作可达到的最优精度；$20$ 米是实际清除门槛，不是所有候选点必须通过的硬条件。

\subsection{一阶端点近似指标（仅用于排序）}\label{sec:hat}

把参考源位置暂时放在第一次示向度中心射线上 $G=(r,0)$，对测向方程作小角度线性化，位置误差满足
\begin{equation}
  |\Delta v|\le r\delta,\qquad
  \bigl|\,y\,\Delta u+(r-x)\,\Delta v\,\bigr|\le\bigl((r-x)^2+y^2\bigr)\delta .
\end{equation}
近似误差区域是中心对称平行四边形，其覆盖半径（等于最大顶点距离）为
\begin{equation}\label{eq:rhohat}
  \hat\rho(r;x,y)=\delta\sqrt{r^2+\frac{\bigl((r-x)^2+y^2+r|r-x|\bigr)^2}{y^2}},\qquad y\ne0 .
\end{equation}
固定 $P$ 后，$\hat\rho$ 关于 $r\ge0$ 经数值验证为凸（验证程序在 $r\in[0,1500]$ 网格上的二阶差分无违反，多组 $P$ 复验），故
\begin{equation}
  \hat f_1(P)=\max\bigl\{\hat\rho(5;P),\,\hat\rho(1500;P)\bigr\}.
\end{equation}
\emph{必须说明：}这是中心射线、一阶线性化且未做真实边界裁剪的近似指标；它\emph{不自动是}真实 $f_1$ 的上界或下界；\emph{不能}单独用于证明 20 米覆盖；\emph{不能}作为永久排除候选点的必要条件。程序中对 $y\approx0$ 的奇异性置为 $10^9$（排名最末），真实有界区域的评价不受此影响。$\hat f_1$ 只用于快速排序和产生搜索起点；为防遗漏，同时保留在 $E$ 内分布较广的多样化起点，避免只搜索近似指标最好的一个小区域。

\subsection{具有必要性依据的粗筛下界}\label{sec:pair}

预先从真实 $D$ 内选取有限个见证源位置 $G_i$（包括远端边界点、若干不同距离与方向的点，共 $13\times7=91$ 个）。它们只是构造下界的见证点，不是声称用有限点替代全部 $D$。

固定 $P$，若两个正常可测见证点 $G_i,G_j$（距 $P$ 均大于 $5$ 米）从 $P$ 看的最小圆周夹角不超过 $2\delta$，则存在同一实际读数 $\beta$ 与两者同时相容，于是 $G_i,G_j\in J(P,\beta)$，故
\begin{equation}
  f_1(P)\ge\frac{\|G_i-G_j\|}{2},
  \qquad
  \ell_{\mathrm{pair}}(P)=\max_{\text{满足条件的点对}(i,j)}\frac{\|G_i-G_j\|}{2}.
\end{equation}
小夹角条件用向量点积判断，避免大量 atan2：
\begin{equation}
  (G_i-P)\cdot(G_j-P)\ge\cos(2\delta)\,\|G_i-P\|\,\|G_j-P\|,
\end{equation}
并注意零向量、浮点容差与近距离分支。没有符合条件的点对时取下界 $0$（而不是认为精度为 $0$）。此外，§\ref{sec:subdiv} 中采样可行读数处的内近似覆盖半径同样提供下界，取
$\ell(P)=\max\bigl\{\ell_{\mathrm{pair}}(P),\ \text{已采样可行}\ \beta\ \text{的覆盖半径下界}\bigr\}$。

\emph{排除规则：}先对若干有希望候选点精细评价，得到当前最好候选点的真实评分上界 $U_{\mathrm{best}}$。若 $\ell(P)>U_{\mathrm{best}}$，$P$ 的真实 $f_1$ 严格更大，可安全淘汰；若希望保留比当前最好结果差不超过 $\tau$ 米的候选区域，则 $\ell(P)>U_{\mathrm{best}}+\tau$ 时才排除。这里必须区分：\emph{近似函数}的差只说明排序靠后；\emph{真实下界}已超过阈值，才支持数学意义上的排除。用离散点判断不能自动证明整块网格单元都被排除，本文始终把粗筛结果标为“采样筛选区域”。

\subsection{角度区间自适应细分与上下界}\label{sec:subdiv}

将可能观测角分为区间 $I_j=[\beta_j-\eta_j,\beta_j+\eta_j]$。对整个区间构造外包区域
\begin{equation}
  J_j^{+}=D\cap W(P;\beta_j,\delta+\eta_j),
\end{equation}
它包含区间内每个实际读数对应的定位区域（读数在区间内时其楔形必含于加宽楔形），故其覆盖半径给出该角度区间的上界 $U_j$；在确实可行的具体角度 $\beta$ 上计算内近似覆盖半径，形成下界。维护 $L(P)\le f_1(P)\le U(P)=\max_j U_j$，每次优先细分\emph{上界最大}、仍可能决定最坏结果的角度区间，并检查：
\begin{enumerate}
  \item $U(P)-L(P)\le$ 指定精度（收敛）；
  \item $L(P)>$ 当前淘汰阈值（提前返回“无需继续评价”）；
  \item 达到计算预算（返回尚未收敛的上下界与明确状态，不假装已经求准）。
\end{enumerate}
粗搜索用约 $1\sim2$ 米评价间隙，最终少量候选点用 $0.1$ 米或更小；这些是可配置算法参数，不是题目给定常数。可行读数范围 $\Theta(P)$ 的安全覆盖区间这样得到：$P$ 在 $D$ 内部时取整个圆周；否则方向集合 $\{\arg(G-P):G\in D\}$ 是连续弧，其端点在 $D$ 边界上，而 $P\in E$ 时 $\|P\|\le1005<1500$，外圆无切点，故端点必在四个角点处，取四角方向极值再加 $0.05^\circ$ 数值余量。初始角度箱取 $1^\circ$ 半宽（保证加宽楔形半角 $2^\circ\ll90^\circ$，半平面表示成立）。\emph{本方法不假设“两个观测角端点足以覆盖全部最坏情况”，也无需手工解析全部关键角度。}

\subsection{两阶段选点算法}\label{sec:algo}

\emph{A.} 用 $E$ 的显式圆盘约束生成一批粗候选点（$20\times20$ 规则网格加 $160$ 个均匀随机点，共约 $360$ 个；不做全 $E$ 密集精确扫描）。\emph{B.} 用 $\hat f_1$ 排序，取近似最优 $20$ 个，并以贪婪最远点法另留 $12$ 个多样化起点。\emph{C.} 按近似优劣依次粗评（角度区间细分，预算 $160$ 次/点，间隙 $2$ 米），建立当前最好上界 $U_{\mathrm{best}}$。\emph{D.} 用点对下界数学排除：$\ell(P)>U_{\mathrm{best}}+\tau$ 的点直接淘汰；粗评中 $L(P)>U_{\mathrm{best}}$ 的点提前终止。\emph{E.} 从保留的若干起点出发做带 $E$ 约束的二维模式搜索（坐标方向，步长 $48\to4$ 米减半），\emph{不假设 $f_1$ 关于 $P$ 凸}，也不只搜 $E$ 边界。\emph{F.} 全局状态缓存：每个已评价点缓存其角度箱状态，提高精度时继续细分而非重新计算；缓存键为毫米量化坐标（本搜索中无两个不同评价点落入同一键）。\emph{G.} 只对 $U\le U_{\mathrm{best}}+3$ 米的少量候选点以 $0.1$ 米间隙精评复核。全程输出：已评价候选点数、完整/提前终止评价次数、当前最好上下界、缓存命中数、耗时与收敛状态；所有耗时均为实测，不以外推的“每点耗时”承诺总时间。

\subsection{精度与移动距离的分层选择策略}\label{sec:f2}

主评价 $f_1$（越小越好）与移动成本 $f_2(P)=\|P\|$（到第二站的移动时间为 $\|P\|/5$，速度 $5$ 米/秒）采用分层优化，不引入任意权重：理论上
\begin{equation}
  m=\inf_{P\in E}f_1(P),\qquad
  F_\tau=\{P\in E:\ f_1(P)\le m+\tau\},
\end{equation}
再在 $F_\tau$ 中选 $\|P\|$ 最小者；$\tau$ 为可配置精度让步，默认取 $1$ 米并作参数对比。注意实际算法没有完成全局最优证明，只能得到数值最佳参考值 $\hat m$；“相对于当前数值最佳参考值的优良候选区域”不等于已证明的全局 $F_\tau$。只优化 $\|P\|$ 对应到第二站的时间；若讨论直到清除的总路程，还要加入从 $P$ 到实际观测后包围圆圆心的路程，两者不得混淆。若需要归一化展示，可定义 $Q(P)=1-f_1(P)/\rho(D)$（$\rho(D)\approx750.1$ 米为 $D$ 自身的最小包围圆半径），表示最坏情况下包围半径的相对缩减程度；第一版未将其加入目标。

\subsection{真实运行结果与参数分析}\label{sec:results}

\paragraph{验证套件（全部通过，$51$ 项）。}最小包围圆基础测试：单点、两点、共线三点、等边三角形、已知圆上 $20$ 点、重复点、近共线点、随机 $10$ 点与枚举对拍（半径误差 $<10^{-8}$）；标准候选点回归：$P=(750,600)$ 实测 $f_1\in[59.332,59.374]$ 米、最坏读数 $319.129^\circ$，与既有临时验证程序的结果（下界 $59.337$、上界 $59.421$ 米）一致且区间更窄；角度周期性：$(750,600)$ 与 $(750,-600)$ 对称一致，跨 $0/360^\circ$ 情形（$P=(750,26)$，最坏读数 $355.06^\circ$）正确；近共线情形 $P=(750,0)$：受 $D$ 限制有界，$f_1\in[372.63,372.71]$ 米，与近似公式 $y=0$ 奇异无关；空交（$\beta=90^\circ$ 不可行）、正常观测与 $5$ 米饱和分支（距 $P=5$ 米的点被内区域排除、$15$ 米的点保留）正确区分；上下界随评价精度（$2/0.5/0.1$ 米）单调收紧且 $U-L\le$ 目标 $+0.02$ 米；固定随机种子重复运行结果一致；独立 $1501$ 点密角度采样最大外半径 $59.28$ 米落在细分上下界之间（密采样只是核验，不是连续角域上界的证明）；一阶近似公式与平行四边形覆盖半径逐点吻合，$\hat\rho$ 关于 $r$ 的凸性数值验证无违反；$E$ 的四圆盘刻画与 $28800$ 点暴力最坏距离判据零失配。

\paragraph{两阶段选点结果（标准完整扇区，$E$ 内）。}
\begin{table}[H]\centering\small
\caption{主要候选点与选择策略实测结果（单位：米，最坏读数单位：度）}
\label{tab:q2results}
\begin{tabular}{lccccc}
\toprule
点 & 坐标 $(x,y)$ & 下界 $L$ & 上界 $U$ & 最坏读数 & 距离 $\|P\|$ \\
\midrule
$P_{\mathrm{best}}$（数值最优参考 $\hat m$） & $(818.47,580.15)$ & $55.657$ & $55.729$ & $317.33$ & $1003.23$ \\
$P_{f_2}$（$\tau=1$ 米距离最优） & $(799.22,-589.98)$ & $56.504$ & $56.593$ & $42.34$ & $993.39$ \\
$P_{f_2}$（$\tau=3$ 米） & $(794.47,-580.15)$ & $57.223$ & $57.303$ & $41.69$ & $983.74$ \\
$P_{\mathrm{hat}}$（仅近似指标 $\hat f_1=60.78$） & $(818.47,580.15)$ & $55.657$ & $55.729$ & $317.33$ & $1003.23$ \\
回归参考 $P=(750,600)$ & $(750,600)$ & $59.332$ & $59.374$ & $319.13$ & $960.47$ \\
\bottomrule
\end{tabular}
\end{table}

数值最优参考 $\hat m=55.729$ 米，最优候选 $P_{\mathrm{best}}=(818.47,580.15)$，其真实评分上下界 $[55.657,55.729]$ 米（间隙 $0.073$ 米），最坏读数 $317.33^\circ$。分层选择（$\tau=1$ 米）：在 $F_\tau$ 内移动距离最短的点 $P_{f_2}=(799.22,-589.98)$，移动距离 $993.39$ 米、移动时间 $198.68$ 秒。在默认参数下，仅用近似指标 $\hat f_1$ 选出的点 $P_{\mathrm{hat}}=(818.47,580.15)$（$\hat f_1=60.78$ 米）与两阶段数值最优参考重合；两阶段方法的增益体现在三处：(a) 给出带上下界与收敛状态的评分而非单一近似值；(b) 粗筛与提前淘汰把真实评价调用压缩到 $106$ 次（纯点对下界排除 $5$ 点、细分中提前终止 $26$ 点、缓存命中 $44$ 次），而候选总数 $411$ 个；(c) 在 $F_\tau$ 内做距离分层选择：$\tau=1$ 米时距离 $993.39$ 米，比直接取 $P_{\mathrm{best}}$ 的 $1003.23$ 米少 $9.83$ 米（约 $2.0$ 秒移动时间）。全部候选约 $411$ 个，粗评 $32$ 个，评价调用 $106$ 次；外层搜索实测耗时 $6.1$ 秒（不含 MATLAB 启动）。图~\ref{fig:q2overview}--\ref{fig:q2dense} 给出候选域、最坏定位区域、近似指标热图、收敛曲线、多边形分辨率与密采样核验；插值热力图仅用于展示，不作为严格评分或整体可行性证明。

\paragraph{参数敏感性。}\begin{table}[H]\centering\small
\caption{参数敏感性}
\label{tab:q2sens}
\begin{tabular}{lcc}
\toprule
参数 & 取值 & 结果 \\
\midrule
$\tau$（精度让步，米） & $0$ & $F_\tau$ 内 $1$ 个已精评点，最短距离 $1003.23$ 米（$200.65$ 秒） \\
$\tau$（精度让步，米） & $1$ & $F_\tau$ 内 $23$ 个已精评点，最短距离 $993.39$ 米（$198.68$ 秒） \\
$\tau$（精度让步，米） & $3$ & $F_\tau$ 内 $40$ 个已精评点，最短距离 $983.74$ 米（$196.75$ 秒） \\
精评间隙（米） & $0.5/0.1/0.05$ & 上下界差 $0.0727/0.0727/0.0365$，均收敛 \\
候选数量 & $291/571$ & $\hat m$ 均为 $55.729$ 米（同一 $P_{\mathrm{best}}$） \\
\bottomrule
\end{tabular}
\end{table}
候选数量从 $291$ 增至 $571$ 个时数值最优参考保持 $55.729$ 米（同一点 $P_{\mathrm{best}}$），结论稳定。精评间隙 $0.5/0.1/0.05$ 米时 $P_{\mathrm{best}}$ 的上下界差为 $0.0727/0.0727/0.0365$ 米，均正常收敛（$0.5$ 与 $0.1$ 米两次命中缓存状态，间隙不变是缓存复用而非未收敛）。

\subsection{计算复杂度、实测耗时与局限性}\label{sec:cost}

单次 $\beta$ 评价：半平面裁剪 $O(N_{\mathrm{arc}}+m)$，Welzl 随机增量算法期望 $O(m)$（$m$ 为顶点数，约 $20\sim2100$）；角度细分每次分裂新增 $3$ 次 $\beta$ 评价，$(750,600)$ 精评至 $0.05$ 米共 $158$ 次角度评价、耗时 $0.35$ 秒；最坏型点（$P\in D$ 内部、针状区域贴孔）约需 $200\sim400$ 次。整体：验证约 $3.4$ 秒，外层搜索约 $6.1$ 秒，总耗时 $45.3$ 秒（R2025a，单机实测）。

局限性（如实说明）：(1) $\hat m$ 是数值最佳参考，\emph{不是已证明的全局最优}，模式搜索只保证在访问邻域内改进；(2) 密角度采样与离散见证点均为核验/下界构造，不是连续角域的上界证明；(3) $\hat f_1$ 不是严格必要条件，只能用于排序；(4) 目标圆域裁切情形采用保守 $E$（安全子区域），未求最大安全区域；(5) 上下界基于浮点实现与显式集合包含关系，不是区间算术形式化认证；(6) 未把“角度成功比例”当作概率，所有结论均为最坏情形意义上的保证。

\begin{figure}[H]\centering
\includegraphics[width=0.66\textwidth]{\figpath fig1_overview.png}
\caption{两阶段选点概览：$D$（浅蓝）、$E$ 边界（黑点）、见证点（灰点）、各类候选状态与最终选择（红色五角星 $P_{\mathrm{best}}$、紫色菱形 $P_{f_2}$、绿色方块 $P_{\mathrm{hat}}$）}
\label{fig:q2overview}
\end{figure}
\begin{figure}[H]\centering
\includegraphics[width=0.62\textwidth]{\figpath fig2_worst_region.png}
\caption{$P_{\mathrm{best}}$ 在最坏读数下的定位区域：内近似（浅蓝填充）、外近似边界（黑线）、最小包围圆（红线）与圆心（红叉）}
\label{fig:q2worst}
\end{figure}
\begin{figure}[H]\centering
\includegraphics[width=0.60\textwidth]{\figpath fig3_hat_heatmap.png}
\caption{近似指标 $\hat f_1$ 在 $E$ 上的分布（插值热图仅用于展示，不是严格评分）}
\label{fig:q2hat}
\end{figure}
\begin{figure}[H]\centering
\includegraphics[width=0.62\textwidth]{\figpath fig4_convergence.png}
\caption{角度区间自适应细分收敛曲线（红：上界 $U$，蓝：下界 $L$；粗线 $P_{\mathrm{best}}$，细线 $(750,600)$）}
\label{fig:q2conv}
\end{figure}
\begin{figure}[H]\centering
\includegraphics[width=0.55\textwidth]{\figpath fig5_bounds_vs_N.png}
\caption{上下界随圆弧多边形分辨率 $N$ 收紧（$P=(750,600)$）}
\label{fig:q2N}
\end{figure}
\begin{figure}[H]\centering
\includegraphics[width=0.60\textwidth]{\figpath fig6_dense_check.png}
\caption{独立密角度采样一致性核验（$P=(750,600)$，密采样只是核验，不是连续角域上界的证明）}
\label{fig:q2dense}
\end{figure}
